% concatenated from main.tex
\documentclass[11pt]{amsart}
\usepackage[margin=1in]{geometry}
\usepackage[multiple]{footmisc}
\usepackage[
colorlinks=true,
linkcolor=blue,
citecolor=blue,
urlcolor=blue]{hyperref}
\usepackage{enumitem}
\usepackage{cleveref}

\DeclareFontShape{OMX}{cmex}{m}{n}{<->cmex10}{}

\DeclareFontFamily{OMX}{MnSymbolE}{}

\DeclareFontShape{OMX}{MnSymbolE}{m}{n}{
    <-6>    MnSymbolE5
    <6-7>   MnSymbolE6
    <7-8>   MnSymbolE7
    <8-9>   MnSymbolE8
    <9-10>  MnSymbolE9
    <10-12> MnSymbolE10
    <12->   MnSymbolE12
}{}

\DeclareSymbolFont{MnLargeSymbols}
                  {OMX}{MnSymbolE}{m}{n}

\DeclareMathDelimiter{\llangle}{\mathopen}
    {MnLargeSymbols}{'164}
    {MnLargeSymbols}{'164}

\DeclareMathDelimiter{\rrangle}{\mathclose}
    {MnLargeSymbols}{'171}
    {MnLargeSymbols}{'171}

\newcommand{\R}{\mathbf R}
\newcommand{\sphere}{\mathbf{S}^{n-1}}
\newcommand{\E}{\mathbf E}
\newcommand{\Prob}{\mathbf P}
\newcommand{\Vol}{\operatorname{vol}}
\newcommand{\rad}{\operatorname{rad}_2}
\newcommand{\DualDvoretzky}{k_\ast}
\newcommand{\Cov}{\operatorname{Cov}}
\newcommand{\Var}{\operatorname{Var}}
\newcommand{\Lip}{\operatorname{Lip}}
\newcommand{\tr}{\operatorname{tr}}
\newcommand{\rank}{\operatorname{rank}}
\let\trace\tr
\newcommand{\SL}{\mathrm{SL}}
\newcommand{\cS}{\mathcal{S}}
\newcommand{\cG}{\mathcal{G}}
\newcommand{\GL}{\mathrm{GL}}
\newcommand{\Gr}{\mathrm{Gr}}
\newcommand{\vrad}{\operatorname{vr}}
\let\vr\vrad
\newcommand{\1}{\mathbf 1}
\newcommand{\vertiii}[1]{{\left\vert\kern-0.25ex\left\vert\kern-0.25ex\left\vert #1 
\right\vert\kern-0.25ex\right\vert\kern-0.25ex\right\vert}}
\newcommand{\gauge}[1]{\vertiii{#1}}
\newcommand{\opnorm}[1]{\vertiii{#1}_{\rm op}}
\newcommand{\fronorm}[1]{\vertiii{#1}_{\rm F}}
\newcommand{\e}{\mathrm{e}}
\newcommand{\ud}{\mathrm{d}}
\newcommand{\PoinConst}{C_{\rm P}}
\newcommand{\ie}{\emph{i.e.,}}
\newcommand{\eg}{\emph{e.g.,}}
\newcommand{\Bodies}{\mathcal{K}_0^n}
\newcommand{\Bobkov}{\mathrm{B}}
\newcommand{\cL}{\mathcal{L}}
\let\vol\Vol
\newcommand{\floor}[1]{\left\lfloor #1 \right\rfloor}
\newcommand{\ceil}[1]{\left\lceil #1 \right\rceil}

\theoremstyle{plain}
\newtheorem{theorem}{Theorem}[section]
\newtheorem{lemma}[theorem]{Lemma}
\newtheorem{proposition}[theorem]{Proposition}
\newtheorem{corollary}[theorem]{Corollary}

\title[Mean width and metric entropy of convex bodies]
{Optimal mean width and metric entropy \\ estimates for convex bodies}

\author{Grigoris Paouris}
\thanks{GP gratefully acknowledges support from the NSF under grant DMS-2405441 and from AFOSR under grant FA9550-25-1-0284.}
\address{Department of Mathematics, 
Texas A\&M University, \and\newline\hspace*{\parindent}%
Department of Mathematics, Princeton University}
\email{grigoris@tamu.edu}

\author{Reese Pathak}
\thanks{RP gratefully acknowledges support from the NSF under grant DMS-2503579.}
\address{Department of Statistics, UC Berkeley,\and\newline\hspace*{\parindent}%
School of Oper. Res. \& Inf. Engineering (ORIE), Cornell University}
\email{pathakr@berkeley.edu}

\begin{document}

\begin{abstract}
We show that there exists a constant $C > 0$ such that for any $n \geq 1$ and any convex body $K \subset \R^n$,
\[
1 \leq \inf_{T \in \GL(n)} \, \frac{M^\ast(TK)}{\vr(TK)} 
 \leq C\sqrt{\log(\e n)},
\]
where $M^\ast$ denotes the spherical mean width and $\vr(\cdot)$ denotes the volume radius. The righthand side is attained, up to universal constants, by the crosspolytope and the regular $n$-simplex. Analogously, we show that, up to universal constants, the logarithm of the Euclidean covering number is maximized over convex bodies $K \subset \R^n$ by the simplex and crosspolytope. 
Our proof makes use of Eldan's stochastic localization.
\end{abstract}

\maketitle

\section{Introduction}


For compact sets $K,L\subset \R^{n}$, with $L$ having
nonempty interior, the covering number $N(K,L)$ is the smallest integer \(m\) for which there exist
\(x_{1},\ldots,x_{m}\in K\) such that
$
K\subseteq \bigcup_{i=1}^{m}(x_{i}+L).
$
Covering numbers quantify the compactness and optimal discretizations of a set; see \S\ref{sec:covering-numbers}.

Our first result shows that for a convex body $K \subset \R^n$, there always exists an ellipsoid that can be used to cover $K$ with no more copies of the ellipsoid than are required to optimally cover the crosspolytope $B^n_1$ of the same volume.
\begin{theorem}
\label{thm:covering-numbers}
There exists $c_{\ref{thm:covering-numbers}}, C_{\ref{thm:covering-numbers}} \in (0, \infty)$ such that for any $n \geq 1$ and any convex body $K\subset\R^{n}$, there exists $T \in \GL(n)$ with $\vol_n(TK) = \vol_n(B^n_1)$ such that \begin{equation}
\label{THM1-1}
\log N\left(TK,rB_{2}^{n}\right)
\leq
  C_{\ref{thm:covering-numbers}}\log N\left( B_{1}^{n},c_{\ref{thm:covering-numbers}}rB_{2}^{n}\right), \quad \mbox{for all}~r > 0.
\end{equation}
Above, $B_{2}^{n}$ denotes the Euclidean unit ball.
\end{theorem}
 

A classical result of Sch\"{u}tt \cite{Sch84}
establishes the Euclidean metric entropy of the 
crosspolytope%
\footnote{We refer to the logarithm of the covering number $\log N(K, rB^n_2)$ as the (Euclidean) \emph{metric entropy} of $K$.}%
\textsuperscript{,}%
\footnote{The notation $a \simeq b$ means that there exist universal constants
$c,C>0$ such that
$
c\,a\leq b\leq C\,a.
$
We reserve $c,C,C_{1}, c', C' >0$ for universal constants.}
\begin{equation}
\label{Carsten}
\log N\bigl(B_{1}^{n},rB_{2}^{n}\bigr)
\simeq 
\begin{cases}
n\log\left(1+\frac{1}{r\sqrt{n}}\right),
& 0<r\leq n^{-1/2}, \\
\frac{1}{r^{2}}\log\left(1+nr^{2}\right),
& n^{-1/2}<r\leq c
\end{cases}.
\end{equation}
The same estimates hold, with adjusted universal constants, when the crosspolytope $B_{1}^{n}$ is replaced by the regular $n$-simplex $S_{n}$.

A fundamental parameter associated with a convex body
is its spherical mean width, 
\[
M^\ast(K) = \E_{\theta \sim \sigma_{n-1}} \Big[ \sup_{x \in K} \, \langle x, \theta \rangle \Big]
= \E_{\sigma_{n-1}} \, h_K(\theta).
\]
Above, $\sigma_{n-1}$ denotes the uniform measure on the sphere $\sphere$. The spherical mean width (or its Gaussian counterpart, the Gaussian width) appears naturally in analysis, probability, and geometry; see~\S\ref{sec:mean-width} for additional remarks. A fundamental question is the maximal order of $M^\ast(K)$ over convex bodies $K\subset \R^n$ with fixed volume.

By Urysohn's inequality (see \cite {Sch14}), it is known that
$M^\ast(K) \geq \vr(K)$, with equality for $K = B^n_2$, where $\vr(K) = (\tfrac{\vol_n(K)}{\vol_n(B^n_2)})^{1/n}$ denotes the volume radius of a convex body. 
To obtain upper bounds, one must first take a suitable linear image $TK$ (also referred to as a \emph{position} of $K$), where $T \in \GL(n)$. The question is then to determine how large the \emph{minimal} mean width,
\[
\inf_{T\in \GL(n)} \frac{ M^\ast(TK) }{ \vr(TK)},
\]
can be over all convex bodies $K \subset \R^n$.
Despite considerable interest, the order of this minimal mean width remained unknown.
Our second result resolves this question: the maximum, up to universal constants, is attained by the $n$-simplex and the crosspolytope.

\begin{theorem}
\label{thm:reverse-urysohn}
There exists $C_{\ref{thm:reverse-urysohn}} \in (0, \infty)$ such that for any $n \geq 1$ and any convex body $K\subset\R^{n}$,
\begin{equation}
\label{THM2-1}
1\leq \inf_{T \in \GL(n)} \, 
\frac{M^\ast(TK)}{{\rm vr}(TK)}
\leq C_{\ref{thm:reverse-urysohn}} \, \frac{M^\ast (B_1^n)}{{\rm vr} (B_1^n)}  \simeq
\sqrt{\log(\e n)}.
\end{equation}
\end{theorem}

The ``Urysohn ratio'' appearing in~\Cref{thm:reverse-urysohn} is the first---and, by the Alexandrov
inequalities, the largest---member of a family of fundamental
geometric parameters associated with a convex body \(K\), namely its
normalized intrinsic volumes. Following the normalization arising from Kubota's
formula~\cite{Sch14}, we set, for an integer $k \in \{1, 2, \dots, n-1\}$,
\begin{equation}
\label{def-W-k}
W_{[k]}(K)
:=
\left(
\int_{\Gr_{n,k}}
\frac{\vol_k(P_F K)}{\vol_k(B_2^k)}
d\nu_{n,k}(F)
\right)^{\frac{1}{k}}.
\end{equation}
We also set $W_{[n]}(K) = \vr(K)$. 
Above, $\Gr_{n,k}$ denotes the Grassmann manifold of $k$-dimensional
subspaces of $\R^n$, and $\nu_{n,k}$ its Haar probability
measure. We denote by $P_F$ the orthogonal projection onto a subspace $F \subset \R^n$. 
Note that by the Alexandrov inequalities, for any convex body $K \subset \R^n$,
\begin{equation}
W_{[1]}(K)\geq W_{[2]}(K)\geq\cdots
\geq W_{[n]}(K).
\end{equation}

Our third result states that there exists a linear transformation such that \emph{all} these quantities are \emph{simultaneously} controlled by
the corresponding ratios for the crosspolytope (or simplex).

\begin{theorem}
\label{thm:quermassintegrals}
There exists $C_{\ref{thm:quermassintegrals}} \in (0, \infty)$ such that for any $n \geq 1$ and every convex body $K \subset \R^n$, there exists $T \in \GL(n)$ such that for $1 \leq k \leq n$,
\begin{equation}
\label{THM3-1}
1 \leq \frac{W_{[k]}(TK)}{{\rm vr}(TK)} \leq C_{\ref{thm:quermassintegrals}} \,
\frac{W_{[k]}(B_1^n)}{ {\rm vr} (B_1^n)} \simeq  \sqrt{\log\left(\frac{\e n}{k}\right)}.
\end{equation}
\end{theorem}
We may also replace the crosspolytope by the $n$-simplex $S_n$ in~\Cref{thm:reverse-urysohn,thm:quermassintegrals}, as 
\[
\frac{W_{[k]}(S_n)}{ {\rm vr}(S_n)}
\simeq \sqrt{\log \Big(\frac{\e n}{k}\Big)}, \quad \mbox{for}~1 \leq k \leq n.
\] 
\subsection{History and additional remarks} 

\subsubsection{Covering numbers}
\label{sec:covering-numbers}
For symmetric convex bodies, covering numbers are precisely the geometric counterpart of entropy numbers of operators between Banach spaces, which are fundamental quantities in operator and approximation theory; see \cite{Pin85, CarSte90} and the references therein. For instance, Sch\"utt's result~\eqref{Carsten} is a direct translation from the original formulation for entropy numbers. 

A breakthrough by V.\ Milman \cite{Mil86} was the realization that by fixing the Euclidean structure in dimension $n$, the entropy numbers for all convex bodies are the same. Milman's M-ellipsoid theorem asserts that for every symmetric convex body $K\subset \R^n$ there exists an ellipsoid $E$ of equal volume, with $\log{N(K, E)} \leq c n$~\cite{Mil86}. This also extends to the non-symmetric case; see \cite{MilPaj00}. \Cref{thm:covering-numbers}, with $r= \vr(B_{1}^{n})$, corresponds to V.\ Milman's relation: \ie~the position used in our result is an M-position (in the sense of V.\ Milman).

Milman's position is a central tool in the local theory of Banach spaces and in asymptotic convex geometry. Among its principal consequences are Milman's reverse Brunn--Minkowski inequality and the Bourgain--Milman ``reverse Santal\'o'' inequality \cite{BouMil87}. Pisier considerably strengthened Milman's theorem by proving the existence of regular ellipsoids. More precisely, for every $\alpha \in (0, 2)$ and every symmetric convex body $K$, there is an ellipsoid $E_\alpha$ such that
\begin{equation}
\label{Pisier}
 \max\left\{ \log{N (K, t E_{\alpha}) },\log{ N( K^{\circ}, t E_{\alpha}^\circ )}\right\} \leq C_{\alpha} \frac{n}{ t^{\alpha}} , \ t >0.
 \end{equation}
See \cite[Chapter 7, Theorem 7.13]{Pis89}. This is known as Pisier's $\alpha$-regular M-position. In the above setting, as $\alpha$ tends to $2$, $C_{\alpha}$ tends to infinity. Thus, \Cref{thm:covering-numbers} introduces the necessary logarithmic correction in $t$ when only one of the two covering numbers is considered. The non-symmetric analogue of \eqref{Pisier} is still an open problem~\cite{Vri24,BizKla26}.

 Artstein, Milman, and Szarek settled the Hilbertian case of Pietsch's long-standing duality of entropy conjecture \cite[Theorem 1]{ArtMilSza04}. Their result implies that when $K$ is symmetric and $E$ is an ellipsoid, the covering numbers satisfy $\log{N(K,E)} \simeq \log{N(E^{\circ},cK^{\circ})}$. Hence,~\Cref{thm:covering-numbers} has an immediate dual counterpart.
 
 Finally, we mention K.\ Ball's observation~\cite{PHDBall} that, assuming the hyperplane conjecture, the isotropic position of a convex body is an M-position. The recent breakthrough paper of Bizeul--Klartag~\cite{BizKla26} confirms that isotropic position is, up to polylogarithmic factors, a $2$-regular Pisier M-position.
 
 \subsubsection{Mean width of convex bodies}
 \label{sec:mean-width}

 The mean width of a convex body is a fundamental geometric parameter with applications in asymptotic convex geometry, the local theory of Banach spaces, the geometry of numbers, high-dimensional probability, random matrix theory, empirical-process theory, compressed sensing, high-dimensional statistics, and convex inverse problems; see, for example, \cite{AubSza17, Ver18, Reg22}.
 
 In probability theory, it frequently appears in its Gaussian normalization: the Gaussian mean width of a bounded set $K$ is $w(K) =\E\sup_{x\in K}\langle G,x\rangle
= \E h_K(G)$. Note that $\sqrt{n} \, M^\ast(K) \simeq w(K)$.
Hence, these parameters are equivalent to the supremum of a canonical Gaussian process. Beginning with the work of Dudley, Sudakov, and Fernique, a close connection between expected suprema of processes and the metric properties (in particular, the metric entropy) of the index set was established. This culminated in Talagrand's majorizing measures theorem \cite{Tal87,Tal05}: the Gaussian width is entirely determined, up to universal constants, by the metric structure of the process.

A parallel development took place in the local theory of Banach spaces, initiated primarily by Maurey and Pisier. For a Banach space $X$, its $K$-convexity constant is the norm $K(X)=\|{\rm {Rad}}_{X}\|$ of the Rademacher projection (see the survey \cite{Mau03}). Pisier's fundamental theorem (see \cite{Pis82}) states that bounded $K$-convexity is equivalent to the absence of uniformly embedded copies of $\ell _1^m$. Figiel and Tomczak--Jaegermann introduced the $\ell$-norm and established its fundamental connection with mean width and $K$-convexity \cite{FigTomJae79}. Their result implies that, for every symmetric convex body $K\subset\R^n$, there exists $T\in \GL(n)$ such that
\begin{equation}
\label{FT}
M(TK) \, M^\ast(TK)\leq CK(X_K),
\end{equation}
where $M(K)=M^\ast(K^{\circ})$ and $X_{K}$ is the Banach space whose unit ball is $K$. Pisier subsequently proved the quantitative estimate
\[
K(X)\leq C\bigl(1+\log d(X,\ell_2^n)\bigr)\leq C\, \log(\e n)
\]
for every $n$-dimensional normed space $X$ \cite{Pis81}; see also \cite[Chapter~2]{Pis89}. Combining these results gives the classical ``$MM^\ast$''-estimate
\[
\inf_{T\in \GL(n)}\, M(TK)\, M^\ast(TK)\leq C\log(\e n).
\]
This yields a (suboptimal) estimate of order $\log n$ for the reversal of Urysohn's inequality.
Note that it extends easily to the non-symmetric case via the Rogers--Shephard inequality.

For Banach lattices, Pisier obtained the sharper estimate
$K(X)\leq C\sqrt{\log(\e n)}$, which leads to optimal estimates in that case. Other notable cases include zonoids and the polars of unit balls of subspaces of $L_{p}$, for which sharp results were previously known \cite{GiaMilRud00,PaoVal18}.

It has been conjectured that every symmetric convex body satisfies
$$
\inf_{T\in \GL(n)} \, M(TK)M^\ast(TK)\leq C\sqrt{\log(\e n)};$$
see the discussion in \cite[p.~207]{AubSza17}. The authors also conjectured a $\sqrt{\log (\e n)}$ upper bound for the normalized minimal mean width, which we establish in~\Cref{thm:reverse-urysohn}.

Prior to our work, the paper by Bizeul--Klartag~\cite{BizKla26} reduced the minimal mean width conjecture to the KLS conjecture. Their work and ours depend on Eldan's stochastic localization~\cite{Eld12, Eld13} in the form introduced by Lee and Vempala~\cite{LeeVem24}. Moreover, these works make use of the idea initiated by Eldan--Lehec~\cite{EldLeh14} to obtain comparison theorems for log-concave isotropic measures via stochastic localization.

\subsubsection{Quermassintegrals and intrinsic volumes of convex bodies}
\label{sec:quermassintegrals}

Let $K$ be a convex body in $\R^{n}$. The Steiner formula states that
\[
\vol_n(K+tB_2^n)
=\sum_{k=0}^{n} \vol_k(B_{2}^{k}) \, V_{n-k}(K) \, t^{k},
\]
where $V_j(K)$ denotes the $j$th intrinsic volume of $K$. In our notation,
\[
V_{k}(K) = \binom{n}{k} \frac{\vol_n(B_{2}^{n})}{ \vol_{n-k}(B_{2}^{n-k})} W_{[k]}^{k}(K). 
\]
The Alexandrov inequalities imply that among convex bodies of fixed volume, every $V_k$ (or $W_{[k]}$) is minimized by Euclidean balls. The case $k=n-1$ corresponds to the classical isoperimetric inequality for surface area, while $k=1$ corresponds to Urysohn's inequality~\cite[Theorem~7.3.1]{Sch14}.

The reverse problem is to minimize $W_{[k]}(TK)$ over volume-preserving linear (or affine) transformations $T$ and determine the bodies for which this affine minimum is largest. A complete sharp answer is presently known only for $k=n-1$. Ball's reverse isoperimetric theorem shows that the extremal body is the cube in the centrally symmetric class and the simplex in the general class \cite{Ball91}. Ball's argument is based on the geometric Brascamp--Lieb inequality; see also Barthe~\cite{Bar98}, who developed the reverse Brascamp--Lieb inequality and treated its equality cases. For $1\leq k\leq n-2$, no comparable sharp identification of the extremal body is known in full generality. \Cref{thm:quermassintegrals} shows that, up to universal constants, the simplex is a maximizer.

Quermassintegrals are also closely connected with metric entropy. Our normalization is intended to make this connection transparent. In particular, it is not difficult to see that $K$ is in Milman's position if and only if
$
W_{[ n/2]}(K)\simeq {\rm vr}(K)
$; see \cite[Chapter 8]{ArtAviGiaMil15}.
  
\subsubsection{Positions of convex bodies}
\label{sec:positions}
For the reverse isoperimetric inequalities mentioned above, one wants to find the linear transformation that minimizes the particular quantity associated with the convex body of interest. This leads to the notion of ``positions'' of convex bodies. As explained in \cite{GiaMil00}, building on results dating back to F.\ John in 1948, these positions typically lead to ``isotropic'' conditions. In \cite{GiaMil00}, the authors provide conditions for minimizing the $k$th quermassintegral. Notably, we \emph{are not} working with these positions, and in fact, they can differ. The position we use was introduced by Bobkov \cite{Bob11} and was shown to be a Milman position. It is known that~\Cref{thm:covering-numbers,thm:quermassintegrals} do \emph{not} hold in the minimal mean width position. Markessinis--Saroglou--Paouris \cite{MarPaoSar12} provide an example of a convex body that is in the minimal mean width position but is far from Milman's position.

\subsection{Ideas in the proof}

We make use of an idea originating with Eldan--Lehec~\cite{EldLeh14} and later presented in an optimized form by Bizeul--Klartag~\cite{BizKla26}. The idea is to use the stochastic localization process to obtain comparison theorems for isotropic log-concave random vectors.

Let $\psi_n$ denote the KLS constant in dimension $n$. The main result in~\cite{BizKla26} shows that, by working with the uniform measure on the convex body $K^\circ$, one obtains
\begin{equation}
\label{ineq:comparison-via-KLS}
\sqrt{n} \, M^\ast(K) \leq C \, \psi_n \, \sqrt{\log (\e n)}.
\end{equation}
Thus, if the KLS conjecture were resolved (\ie~if $\psi_n = O(1)$), then it would imply~\Cref{thm:reverse-urysohn}.

Thus, we need to avoid the reduction to the KLS conjecture. Our key idea in this direction is to replace the uniform measure on $K^\circ$ with the Gaussian measure conditioned on $K^\circ$. Assuming that this measure has scalar covariance is equivalent to assuming that $K^\circ$ lies in Bobkov's maximal Gaussian measure position~\cite{Bob11}. The measures arising along the stochastic localization process also immediately enjoy bounded Poincar\'{e} constants. By controlling the third-order tensor arising in the analysis of the stochastic localization process in terms of the Poincar\'e constant (\eg~\Cref{lem:mean-zero-log-ccv}), we obtain the sharp reversal of Urysohn's inequality, \Cref{thm:reverse-urysohn}.

To obtain the bounds on the quermassintegrals (\Cref{thm:quermassintegrals}), the covering numbers (\Cref{thm:comparison-theorem}) and a sharpened version of~\Cref{thm:reverse-urysohn} involving the Dvoretzky number (see~\Cref{thm:main}), we require another idea. Observe that the factor $\sqrt{\log n}$ which appears in inequality~\eqref{ineq:comparison-via-KLS} is due to controlling the time at which the minimal eigenvalue $\lambda_n(A_t)$, with $A_t = \Cov(\mu_t)$ denoting the covariance associated with the stochastic localization process $\mu_t$, drops below a universal constant. We prove a more general comparison inequality (see~\Cref{thm:comparison-theorem} for a rigorous statement),
\begin{equation}
    \label{ineq:comparison-ineq-intro}
\E \gauge{G} \leq C \,\min_{1 \leq m \leq n } \, \Big(\sqrt{\log(\e n/m)} \, \E \gauge{X} + \Lip(\gauge{\cdot})\sqrt{m} \Big),
\end{equation}
when $X$ is distributed according to a conditioned Gaussian measure in Bobkov's position and $\gauge{\cdot}$ is a gauge on $\R^n$. This inequality is established by controlling the stopping time for the $m$th smallest eigenvalue $\lambda_{n-m+1}(A_t)$, \ie~the first time at which this eigenvalue drops below a universal constant. The factor $\log(\e n/m)$ arises from the fact that the stopping time for the eigenvalue $\lambda_{n-m+1}(A_t)$ can be bounded on the order of $\log^{-1}(\e n/m)$. We obtain our covering number bounds by combining the comparison inequality~\eqref{ineq:comparison-ineq-intro} with the well-known relations between the negative moments of the spherical mean width, the quermassintegrals, and the width and covering numbers of the intersections $K \cap rB^n_2$, in a form given by Mourtada~\cite{Mou25}.

\section{Preliminaries}

A convex body is a compact, convex set with nonempty interior. We denote by $\Bodies$ the class of all centrally symmetric convex bodies $K \subset \R^n$.


\subsection{\texorpdfstring{$\Bobkov$}{B}-position and conditioned Gaussian measures}

We will make use of conditioned Gaussian measures on $\R^n$. We denote the standard Gaussian measure and the Gaussian measure with variance $\tau^{-2}$ by
\[
\gamma_n = N(0, I_n), \quad \mbox{and} \quad 
\gamma_{\tau, n} = N(0, \tau^{-2} I_n).
\]
The conditioned versions of these measures on some measurable set $L \subset \R^n$ are given by 
\[
\gamma_{L}(A) = \frac{\gamma_n(A \cap L)}{\gamma_n(L)} 
\quad \mbox{and} \quad 
\gamma_{\tau, L}(A) = \frac{\gamma_{\tau,n}(A \cap L)}{\gamma_{\tau, n}(L)} 
\]
for every measurable set $A \subset \R^n$ and $\tau > 0$. 

\medskip 

Following Bobkov~\cite{Bob11}, we say that $K \in \Bodies$ is in \emph{Bobkov's maximal Gaussian measure position} (abbreviated \emph{$\Bobkov$-position}) if
\begin{equation}\label{eqn:variational-condition}
\gamma_n(K) \geq \gamma_n(TK) 
\qquad\text{for every} \quad 
T\in\SL(n).
\end{equation}
Using the first-variation condition for the maximality of $T = I_n$
in the definition of $\Bobkov$-position (see~\eqref{eqn:variational-condition}) together with V.\ Milman's theory of $M$-ellipsoids~\cite{Mil86}, Bobkov~\cite{Bob11} deduces the following properties of $\Bobkov$-position.

\begin{proposition}
[{\cite[Proposition 3.11, p.~32]{Bob11}}]
\label{prop:bobkov}
    If $K \in \Bodies$ is
	    in the $\Bobkov$-position, then the following statements hold:
    \begin{enumerate}[label=(\roman*)]
    \item 
    \label{item:scalar-covariance}
    the covariance matrix of the conditioned measure is scalar, \ie 
\[
\Cov(\gamma_K) = \alpha_{K}^2 I_n;
\]    
    \item
    \label{item:bound-on-alpha}
     if $\vol_n(K) = 1$, then there exists a universal constant $c_{\ref{prop:bobkov}} > 0$ such that 
     $\alpha_{K}^2 \geq c_{\ref{prop:bobkov}}$. 
    \end{enumerate}
\end{proposition}


We also need the following bound in $\Bobkov$-position.
\begin{lemma}
\label{lem:bounded-Lipschitz-constant}
There exists $c_{\ref{lem:bounded-Lipschitz-constant}} > 0$ such that for any $n \geq 1$ and any $K \in \Bodies$ with $K^\circ$ in $\Bobkov$-position and $\vol_n(K^\circ) = 1$, $\rad(K) \leq c_{\ref{lem:bounded-Lipschitz-constant}}$.
\end{lemma}
\begin{proof}
Pick $y \in K$ such that $\|y\|_2 = \rad(K)$. Let $\theta = \tfrac{y}{\|y\|_2}$. Then every $x \in K^\circ$ satisfies $|\langle x, \theta \rangle | \leq \tfrac{1}{\rad(K)}$.
Hence, by \Cref{prop:bobkov}, since $\Cov(\gamma_{K^\circ}) = \alpha_{K^\circ}^2 I_n$, the lower bound on $\alpha_{K^\circ}$ gives
\[
c_{\ref{prop:bobkov}} \leq \alpha_{K^\circ}^2 = \Var_{\gamma_{K^\circ}}(\langle X, \theta\rangle) \leq \frac{1}{\rad(K)^2}.
\]
Rearranging gives $\rad(K) \leq 1/\sqrt{c_{\ref{prop:bobkov}}}$, as claimed.
\end{proof}




We also recall the Poincar\'{e} constant of a probability measure $\nu$ on $\R^n$. It is the smallest $C > 0$ such that 
\[
\Var_\nu(f) \leq C\, \E_\nu \|\nabla f\|_2^2 
\]
holds for all locally Lipschitz $f \colon \R^n \to \R$. We denote the smallest such constant by $\PoinConst(\nu)$. 


\begin{lemma}
\label{lem:poin-const-of-conditioned-gaussians}
Let $\nu = \cL(Z \mid \{Z \in K\})$, where $K \in \Bodies$ and $Z \sim N(z, \tau^2 I_n)$ for some $\tau > 0$ and some $z  \in \R^n$. Then $\PoinConst(\nu) \leq \tau^2$. 
\end{lemma}
\begin{proof}
Note that the Lebesgue density of $\nu$ satisfies
\[
\frac{\ud \nu}{\ud x}(y) = 
\mathcal{Z}^{-1}\exp(-\Phi(y)) \quad
\mbox{where}
\quad \Phi(y) = \begin{cases} 
\|z-y\|_2^2/(2\tau^2), & y \in K \\ 
+\infty,& y \not \in K. 
\end{cases}.
\]
Above, $\mathcal{Z} = \int_K \exp(-\Phi(y)) \, \ud y$. Now observe that
\[
\nabla^2 \Phi(y) = \frac{1}{\tau^2} I_n \quad \mbox{for}~y \in \mathrm{int}(K).
\]
Hence, following the argument on p.~37 of~\cite{BarCorEra13}, we obtain 
\[
\Var_\nu(f) \leq \tau^2 \int_{\mathrm{int}(K)} \|\nabla f(x)\|_2^2 \, \ud \nu(x) \leq \tau^2 \E_{X \sim \nu} \|\nabla f(X)\|_2^2,
\]
which shows that $\PoinConst(\nu) \leq \tau^2$, as needed.
\end{proof}


\subsection{Background on stochastic localization}

We make use of Eldan's stochastic localization process~\cite{Eld12, Eld13}, which we introduce now. Let $\mu$ be a probability measure on $\R^n$ which has compact support and is absolutely continuous with respect to Lebesgue measure. We write $p = \tfrac{\ud \mu}{\ud x}$ for its density. 
For $t \geq 0$ and $\theta \in \R^n$, we define the probability measure $\mu_{t, \theta}$ which satisfies
\[
\frac{\ud \mu_{t,\theta}}{\ud \mu}(x) = \e^{\langle \theta, x \rangle - t \|x\|_2^2/2 - 
\Lambda_t(\theta)}, 
\quad 
\Lambda_t(\theta) = 
\log \E_{X \sim \mu} \e^{\langle \theta, X\rangle - t \|X\|_2^2/2}.
\]
We make use of the following notation for the mean and covariance:
\[
a_t(\theta) = \E_{X \sim \mu_{t, \theta}}[X] 
\quad 
\mbox{and} \quad 
A_t(\theta) = \Cov(\mu_{t, \theta}).
\]
Let $\{B_t\}_{t \geq 0}$ be a standard Brownian motion in $\R^n$, with initial condition $B_0 = 0$. Consider the process $\{\theta_t\}_{t \geq 0}$ defined as the solution of the stochastic differential equation:
\[
\ud \theta_t = a_t(\theta_t)\, \ud t + \, \ud B_t,
\]
with initial condition $\theta_0 = 0$.
With a slight abuse of notation, we refer to the 
family of (random) probability measures $\mu_t \equiv \mu_{t, \theta_t}$ for $t \geq 0$ as the \emph{stochastic localization process}. We also write 
\[
a_t \equiv a_t(\theta_t), \quad A_t \equiv A_t(\theta_t), \quad t \geq 0,
\]
for the corresponding (random) mean and covariance processes. 
We refer to a probability measure $\mu$ as being isotropic if, for $X \sim\mu$,\footnote{To avoid confusion, a random vector $X \sim \mu$ (or its law $\mu$) has \emph{scalar covariance} if $\Cov(\mu)=\Cov(X) = \rho I$ for some $\rho > 0$. We reserve \emph{isotropic} for the case $\rho = 1$.}
\[
\E_{\mu} X = 0, \quad \mbox{and} \quad 
\E_{\mu}\, \Big[X \otimes X\Big] =  I_n.
\]

\section{Comparison theorem for conditioned Gaussian measures}

The main result in this section is the following comparison inequality between the Gaussian measure $\gamma_n$ and the conditioned Gaussian measure $\gamma_{K^\circ}$ for $K^\circ$ in Bobkov's position. Recall that a \emph{gauge} is a nonnegative, positively $1$-homogeneous, convex function.


\begin{theorem}
\label{thm:comparison-theorem}
There exists a constant~$C_{\ref{thm:comparison-theorem}} \in (0, \infty)$ such that for any $n \geq 1$ and any $K \in \Bodies$ with $K^\circ$ in $\Bobkov$-position and $\vol_n(K^\circ) = 1$, we have
\[
\E_{\gamma_n} \gauge{G} \leq C_{\ref{thm:comparison-theorem}}
 \Big(\sqrt{\log(\e n/m)} \E_{\gamma_{K^\circ}} \gauge{X}  + \Lip\big(\gauge{\cdot}\big) \sqrt{m} \Big),
\]
for any gauge $\gauge{\cdot} \colon \R^n \to \R_+$ and any $m \in\{1, 2, \dots, n\}$. 
\end{theorem}


The next result shows that control of the lower eigenvalues of $\{A_t\}_{t \geq 0}$ up to time $T \geq 0$ yields a comparison between the expected norm under the Gaussian measure and that under a conditioned Gaussian measure. We use the notation
\[
\lambda_1(A) \geq \lambda_2(A) \geq \cdots \geq \lambda_n(A)
\]
for the (ordered) eigenvalues of a real, $n \times n$ symmetric matrix $A$. 
\begin{proposition}
\label{prop:from-eigen-control-to-comparison-theorem}
Let $X \sim \mu$, where $\mu$ is a centered probability measure on $\R^n$. Suppose that for $T > 0$ and some $m \in \{1, 2, \dots, n\}$, the stochastic localization process $\{\mu_t\}_{t \geq 0}$ satisfies
    \begin{equation}
    \label{eqn:control-of-eigen-val-event}
    \Prob\bigg\{ \inf_{t \in [0, T]}\lambda_{n-m+1}(A_t) \leq \frac{1}{2} \bigg\} \leq \frac{m}{n}. 
    \end{equation}
    Then, for any gauge $\gauge{\cdot}$ on $\R^n$, we have
    \[
    \E_{\gamma_n} \gauge{G} \leq \frac{2}{\sqrt{T}}\,  \E_\mu \gauge{X} + \Lip\big(\gauge{\cdot}\big)\, \sqrt{2m}.
    \]
\end{proposition}

To control the probability in~\eqref{eqn:control-of-eigen-val-event}, it suffices that the family of measures $\{\mu_t\}_{t \leq T}$ has uniformly small Poincar\'{e} constants. This observation, in various forms, appears in many works on stochastic localization (\eg~Eldan~\cite[Section 3]{Eld13}, Klartag--Lehec~\cite[Lemma~5.2]{KlaLeh22}, Bizeul~\cite[Lemma~3.6]{Biz24}, Lee--Vempala~\cite[Lemma 38]{LeeVem24}, Bizeul--Klartag~\cite[Proposition 2.3]{BizKla26}).

\begin{proposition}
\label{prop:eigen-stop-upper}
There exists $C_{\ref{prop:eigen-stop-upper}} > 0$ 
such that the following holds. 
For any positive integers $m \leq n$,
any $\delta > 0$, and any $K \in \Bodies$ for which $\gamma_{\delta, K}$ is isotropic, the inequality~\eqref{eqn:control-of-eigen-val-event} holds with $T = C_{\ref{prop:eigen-stop-upper}} \delta^6/(\log(\e n/m))$. 
\end{proposition}

Note that in~\Cref{prop:eigen-stop-upper}, since $\gamma_{\delta, K}$ is assumed isotropic, testing with linear functionals and applying~\Cref{lem:poin-const-of-conditioned-gaussians} show that $\delta \leq 1$.


\subsection{Proof of~\texorpdfstring{\Cref{prop:from-eigen-control-to-comparison-theorem}}{the comparison from eigenvalue control}}

\label{app:proofs-of-stoch-loc-lemmas}


We will require the following comparison lemma. The proof is a simple conditioning argument and Jensen's inequality, as given below. An alternative formulation, with a proof using the heat semigroup, is given in 
~\cite[Proposition 3.2]{BizKla26}. In the result below, we denote by $\cS^n_+$ the family of real, symmetric, positive semidefinite $n \times n$ matrices.

\begin{lemma}
\label{lem:comparison-result}
Let $\{M_t\}_{0\leq t\leq T}$ be an $\R^n$-valued, continuous $L^2$ martingale, with $M_0=0$. 
Suppose that, for some predictable, $\cS^n_+$-valued process $\{C_t\}_{0 \leq t \leq T}$, $\ud [M]_t = C_t \, \ud t$. If, for some $\rho>0$,
$C_t \succeq \rho I_n$ holds almost surely, then 
for every convex function $F \colon \R^n\to\R$,
\[
\E F(M_T)\geq \E F\bigl(\sqrt{\rho}\,B_T\bigr),
\]
where $\{B_t\}_{0 \leq t\leq T}$ is a standard Brownian motion in $\R^n$ indexed by $[0, T]$.
\end{lemma}

\begin{proof}
We work on a probability space enlarged to contain an independent standard Brownian motion $\{W_t\}_{0 \leq t \leq T}$ in $\R^n$, indexed by $[0, T]$. Define the continuous local martingale
\[
\ud X_t
= \sqrt{\rho}\,C_t^{-1}\, \ud M_t
+
 \Big(I_n-\rho C_t^{-1}\Big)^{1/2} \, \ud W_t.
\]
The quadratic variation satisfies
\[
\ud [X]_t
=
\rho C_t^{-1} \ud [M]_t C_t^{-1}
+
\Big(I_n-\rho C_t^{-1}\Big) \, \ud t 
=
I_n \, \ud t.
\]
Thus, $\{X_t\}_{0 \leq t \leq T}$ is a standard Brownian motion on $[0, T]$; in particular, it is also a continuous martingale.
Furthermore, the quadratic covariation between $\{M_t\}_{t \leq T}$ and $\{X_t\}_{t \leq T}$ satisfies
\[
\ud[M,X]_t
=
C_t\Big(\sqrt{\rho}\,C_t^{-1}\Big)^{\mathsf T} 
\ud t 
=
\sqrt{\rho}\,I_n \ud t.
\]
Therefore, we can write $M_t = N_t + \sqrt{\rho} X_t$, 
for some continuous martingale $\{N_t\}_{0 \leq t \leq T}$ with $[N,X]_t=0$. 
Now, consider the $\sigma$-field
\[
\cG_t = \sigma(X_s:0\leq s\leq t).
\]
We claim that $\E[N_T \mid  \cG_T]=0$. 
To prove this, let $Z\in L^2(\cG_T)$ be an arbitrary scalar random variable. By the martingale
representation theorem, there exists a $\cG_t$-predictable, $\R^n$-valued process $h_t = (h_t^1, \dots, h_t^n)$ such that
\[
Z=\E Z+ \sum_{j=1}^n \int_0^T h_s^j \,dX_s^j.
\]
Hence, for $i \in \{1, 2, \dots, n\}$, using $\E N_T^i = 0$,  
\[
\E[N_T^i Z] = \E[N_t^i (Z-\E Z)]
=
\sum_{j=1}^n \E \int_0^T h_s^{j}\,d[N^i,X^j]_s
=
0.
\]
Since $Z \in L^2(\cG_T)$ was arbitrary, $\E[N_T \mid \cG_T] = 0$.  Therefore, $\E[M_T\mid\mathcal{G}_T]=\sqrt{\rho}\,X_T$.  
This implies, by Jensen's inequality, that 
\[
F(\sqrt{\rho} X_T) = F(\E[M_T \mid \cG_T]) \leq \E[F(M_T) \mid \cG_T],
\]
almost surely. Passing to the full expectation yields the claim.
\end{proof}

Let $\phi(u) = \max\{u, 1/2\}$. We write the eigendecompositions 
\[
A_t = \sum_{i=1}^n \lambda_i^{(t)} u_i^{(t)} \otimes u_i^{(t)}, \quad \mbox{and} \quad 
\Sigma_t 
= \sum_{i=1}^n \phi\big(\lambda_i^{(t)}\big) u_i^{(t)} \otimes u_i^{(t)},
\]
where $u_i^{(t)}$ and $\lambda_i^{(t)}$ denote eigenvector-eigenvalue pairs of the covariance matrix $A_t$.
We define 
\begin{subequations}
\begin{equation}
\label{eqn:martingale-truncated}    
v_t = \int_0^t \, \Sigma_s \, \ud B_s,
\end{equation}
which is a continuous martingale. Note that 
it also holds that 
\begin{equation}
\label{eqn:martingale-not-truncated}    
a_t = \int_0^t A_s \, \ud B_s.
\end{equation}
\end{subequations}
Consider the stopping time 
\[
\tau = \inf\{t \geq 0 : \lambda_{n - m + 1}(A_t) < 1/2 \}. 
\]
For any $t \geq 0$, it holds that
\begin{equation}
\label{ineq:operator-norm-bound-on-gap}
\opnorm{\Sigma_t - A_t} \leq \frac{1}{2}.
\end{equation}
Hence, for $t \leq \tau$, it evidently holds that 
\begin{equation}
\label{ineq:rank-bound-on-gap}
\rank(\Sigma_t - A_t) \leq m, 
\quad \mbox{and} \quad
\fronorm{\Sigma_t - A_t} \leq \frac{\sqrt{m}}{2}.
\end{equation}
By It\^o's isometry,
\begin{align}
\E \|v_T - a_T\|_2^2 &= 
\E \int_0^{T \wedge \tau} \fronorm{\Sigma_t - A_t}^2 
\, \ud t + 
\E \int_{T \wedge \tau}^T\fronorm{\Sigma_t - A_t}^2 \, \ud t \nonumber \\ 
&\leq \frac{m T}{4} + \E \int_0^T \1\{\tau \leq t\} \fronorm{\Sigma_t - A_t}^2 \, \ud t 
\nonumber  \\ 
&\leq \frac{mT}{4} + 
\frac{n T}{4} \Prob\{\tau \leq T\} 
\nonumber \\ 
&\leq \frac{m T}{2}.
\label{ineq:final-bound-on-gap}
\end{align}
Above, we applied
the bounds in~\eqref{ineq:operator-norm-bound-on-gap}
and~\eqref{ineq:rank-bound-on-gap}, and then the assumption
in~\eqref{eqn:control-of-eigen-val-event}.
Hence, the triangle inequality and inequality~\eqref{ineq:final-bound-on-gap} yield
\begin{equation}
\label{ineq:bound-on-v-t}
\E \gauge{v_T}\leq \E \gauge{a_T} + \E \gauge{v_T -a_T} 
\leq \E \gauge{a_T} + \Lip\big(\gauge{\cdot} \big)\frac{\sqrt{m T}}{\sqrt{2}}.
\end{equation}
Now, we can apply~\Cref{lem:comparison-result} to $F = \gauge{\cdot}$ and $M_t = v_t$. Note that $\ud [v]_t = \Sigma_t^2 \, \ud t$ and $\Sigma_t^2 \succeq \tfrac{1}{4} I_n$, so we may take $\rho = 1/4$. Using~\eqref{ineq:bound-on-v-t}, we obtain
\begin{equation*}
\frac{1}{2} \sqrt{T} \E_{\gamma_n} \gauge{G} 
= \frac{1}{2} \E \gauge{B_T} \leq \E \gauge{v_T} \leq 
\E \gauge{a_T} + \Lip\big(\gauge{\cdot} \big)\frac{\sqrt{m T}}{\sqrt{2}}
\leq 
\E_\mu \gauge{X} + \Lip\big(\gauge{\cdot} \big)\frac{\sqrt{m T}}{\sqrt{2}}.
\end{equation*}
The last inequality follows from Jensen's inequality (\eg~\cite[ineq.~(55)]{BizKla26}).
Rearranging yields the claimed result.

\subsection{Proof of~\texorpdfstring{\Cref{prop:eigen-stop-upper}}{the eigenvalue stopping-time estimate}}
\label{sec:proof-eigen-stop-upper}

The next lemma bounds the Poincar\'{e} constants of the measures arising from stochastic localization on 
$\gamma_{\delta, K}$. 

\begin{lemma}
\label{lem:computation-for-conditional-measure}
Let $\mu = \gamma_{\delta, K}$, where $\delta > 0$ and $K \in \Bodies$. With probability $1$, the stochastic localization process satisfies
\[
\opnorm{A_t} = \opnorm{\Cov(\mu_t)} 
\leq \frac{1}{\delta^2 + t} \leq \frac{1}{\delta^2}
\quad \mbox{for all}~t \geq 0.
\]    
\end{lemma}
\begin{proof}
At time $t$, the measure $\mu_t$ has density relative to $\mu$ given by
\[
\frac{\ud \mu_t}{\ud \mu}(x) \propto \e^{\langle \theta_t, x\rangle - t\|x\|_2^2/2}.
\]
In particular, since $\mu = \gamma_{\delta, K}$, 
completing the square shows that for $X_t \sim \mu_t$,
\[
\cL(X_t) = \cL(Z_t \mid \{ Z_t \in  K \})
\quad \mbox{where} \quad 
Z_t \sim N\Big(\frac{\theta_t}{\delta^2 + t}, \frac{1}{\delta^2 + t} I_n\Big).
\]
By~\Cref{lem:poin-const-of-conditioned-gaussians}, for all $t \geq 0$, we have
\[
\PoinConst(\mu_t) \leq \frac{1}{\delta^2 + t} \leq \frac{1}{\delta^2}. 
\]    
Hence, by considering linear functionals, $\opnorm{A_t} \leq \PoinConst(\mu_t) \leq 1/\delta^2$. 
\end{proof}


We will need the following lemma, which controls the third-order tensors of a log-concave random vector. 
\begin{lemma}
\label{lem:mean-zero-log-ccv}
Let $Z \sim \mu$ be a mean-zero random vector in $\R^n$. Define the tensor 
\begin{equation}\label{eqn:three-tensor}
H_{\theta} = \E\Big[\langle Z, \theta \rangle\,  Z \otimes Z \Big], \quad \theta \in \sphere.
\end{equation}
It follows that
\[
\opnorm{\sum_{i=1}^n H_{e_i}^2} \leq 4\, \PoinConst(\mu) \, \opnorm{\Cov(\mu)}^2.
\]
\end{lemma}
\begin{proof}
Note that $H_{e_i} \theta = H_\theta e_i$ for all $\theta \in \sphere$ and all $i \in \{1, 2, \dots, n\}$. Therefore, 
\[
\opnorm{\sum_{i=1}^n H_{e_i}^2} = \sup_{\theta \in \sphere} \sum_{i=1}^n \langle H_\theta^2 e_i, e_i\rangle = \sup_{\theta \in \sphere}\fronorm{H_\theta}^2. 
\] 
For $\theta \in \sphere$, since $Z$ is centered, we have
\begin{align*}
\fronorm{H_\theta}^4
= \E[\langle Z, \theta \rangle \langle H_\theta Z, Z \rangle]^2 = 
\E[\langle Z, \theta \rangle (\langle H_\theta Z, Z \rangle - \E\langle H_\theta Z, Z \rangle) ]^2.
\end{align*}
Cauchy--Schwarz and the Poincar\'{e} inequality yield the following bound:
\begin{multline*}
\fronorm{H_\theta}^4 
\leq \langle \Cov(Z) \theta, \theta\rangle
\Var_\mu( \langle H_\theta Z, Z \rangle) 
\\ \leq
4 \PoinConst(\mu) \opnorm{\Cov(Z)} \trace(H_\theta^2 \Cov(Z)) 
\leq 4 \PoinConst(\mu) \opnorm{\Cov(Z)}^2 \fronorm{H_\theta}^2. 
\end{multline*}
Rearranging the inequality above furnishes the claim.
\end{proof}

We also require the following consequence of a reverse H\"older inequality for log-concave measures. 

\begin{lemma}[{\cite[Lemma~A.2]{BizKla26}}]
\label{lem:three-tensor-log-ccv}
    Let $Z \sim \mu$ be a mean-zero log-concave random vector in $\R^n$. Then, 
    the tensor $H_\theta$ defined in display~\eqref{eqn:three-tensor} satisfies 
    \[
    \sup_{\theta \in \sphere} \opnorm{H_\theta}^2 \leq 9 \opnorm{\Cov(\mu)}^3. 
    \]
\end{lemma}



We also recall the following maximal inequality for continuous martingales~\cite[Appendix~B]{ShoWel09}.
\begin{lemma}
\label{lem:maximal-inequality-for-cts-martingales}
Let $T > 0$ and let
$\{M_t\}_{t \geq 0}$ be a continuous martingale with $M_0 = 0$. If the quadratic variation satisfies $[M]_T \leq b$ almost surely, then
\[
\Prob\Big\{\sup_{t \in [0, T]} M_t \geq a\Big\} \leq 
\exp\Big\{-\frac{a^2}{2b}\Big\},
\]
for any $a > 0$. 

\end{lemma}


We now proceed to prove~\Cref{prop:eigen-stop-upper}. We use $\llangle X, Y \rrangle = \tr(XY)$ for the trace inner product of real symmetric $n \times n$ matrices $X$ and $Y$. Define the smoothing function by
\[
f_{\beta, m}(A) = \frac{-1}{\beta} \log \tr\Big(\frac{1}{m}\e^{-\beta A}\Big).
\] 
Note that for a real positive semidefinite matrix $A$, we have
\[
f_{\beta, m}(A) \leq \frac{-1}{\beta}\log \frac{1}{m} \sum_{i=1}^m \e^{-\beta \lambda_{n-m+i}(A)} \leq \lambda_{n-m+1}(A). 
\]
At $t = 0$, $A_0 = I_n$, and hence 
\[
f_{\beta, m}(A_0) = 1 - \frac{\log(n/m)}{\beta}.
\]
For $\theta \in \sphere$, define
\[
H_{\theta,t} = \E_{X \sim \mu_t} \Big[\langle X - a_t, \theta\rangle (X-a_t) \otimes (X-a_t)\Big].
\]
A direct computation shows that $\nabla f_{\beta, m}(A_t) \succeq 0$ and $\trace(\nabla f_{\beta, m}(A_t)) = 1$. Hence,
\[
\llangle \nabla f_{\beta, m}(A_t), A_t^2 \rrangle \leq \opnorm{A_t^2} \leq \opnorm{A_t}^2. 
\]
Additionally, by the same computation as in~\cite[Lemma~A.1]{BizKla26},
\[
\sum_{i=1}^n\nabla^2 f_{\beta, m}(A_t)[H_{e_i, t}, H_{e_i, t}] \geq - \beta \Big \llangle \nabla f_{\beta, m}(A_t), \sum_{i=1}^n H_{e_i, t}^2 \rrangle 
\geq - \beta \, \opnorm{\sum_{i=1}^n H_{e_i, t}^2}.
\]
Hence, by It\^o's formula,
\begin{align}
\label{ineq:lower-differential-minus}
\ud f_{\beta,m}(A_t) 
&=\Big\llangle \nabla f_{\beta,m}(A_t), \sum_{i=1}^n H_{e_i, t} \, \ud B_{i,t} \Big\rrangle  - 
\Big\llangle \nabla f_{\beta,m}(A_t), A_t^2 \Big\rrangle \, \ud t + 
\frac{1}{2} \sum_{i=1}^n \nabla^2 f_{\beta, m}(A_t)[H_{e_i, t}, H_{e_i, t}] \, \ud t  
\nonumber \\
&\geq  
\Big\llangle \nabla f_{\beta,m}(A_t), \sum_{i=1}^n H_{e_i, t} \, \ud B_{i,t}\Big\rrangle - \bigg(\opnorm{A_t}^2 + 
\frac{\beta}{2} \opnorm{\sum_{i=1}^n H_{e_i, t}^2}\bigg)\, \ud t.
\end{align}
By~\Cref{lem:computation-for-conditional-measure,lem:mean-zero-log-ccv}, 
\begin{equation*}
\ud f_{\beta,m}(A_t)  
\geq 
\Big\llangle \nabla f_{\beta,m}(A_t), \sum_{i=1}^n H_{e_i, t} \, \ud B_{i,t}\Big\rrangle - \bigg(1 + 
\frac{2\beta}{\delta^2} \bigg)\frac{1}{\delta^4} \, \ud t.
\end{equation*}
Define the continuous martingale 
\[
N_t = \int_0^{t} 
\Big \llangle \nabla f_{\beta,m}(A_s), \sum_{i=1}^n H_{e_i, s} \, \ud B_{i,s}\Big\rrangle.
\]
For $t\leq T$, the quadratic variation of $N_t$ satisfies the following bound:
\[
[N]_t = \int_0^t \sum_{i=1}^n\llangle \nabla f_{\beta,m}(A_s), H_{e_i, s}\rrangle^2 \,\ud s \leq  9 \int_0^{t} \opnorm{A_s}^3 \, \ud s \leq \frac{9}{\delta^6} T. 
\]
Above, we used Cauchy--Schwarz, \Cref{lem:three-tensor-log-ccv}, and the log-concavity of $\gamma_{\delta, K}$ to conclude
\[
\sum_{i=1}^n\llangle \nabla f_{\beta,m}(A_s), H_{e_i, s}\rrangle^2 
= 
\sup_{\theta \in \sphere} \Big \llangle \nabla f_{\beta,m}(A_s),  H_{\theta, s}\Big\rrangle^2
\leq 
\sup_{\theta \in \sphere} \opnorm{H_{\theta, s}}^2 \leq 9 \opnorm{A_s}^3. 
\]
In particular, since $\delta \in (0, 1]$, setting $C = 9$ gives
\begin{equation}
\ud f_{\beta,m}(A_t)  
\geq 
\ud N_t - C \delta^{-6} (1 + \beta)\, \ud t, \quad 
[N]_t \leq C \delta^{-6} t.
\label{ineq:final-diff-lower-minus}
\end{equation}
Integrating~\eqref{ineq:final-diff-lower-minus}, 
for $t \leq T$, we have
\begin{equation*}
f_{\beta,m}(A_{t}) \geq N_t + 1 - \frac{\log(n/m)}{\beta} - C \delta^{-6} T  - C \delta^{-6} \beta T.
\end{equation*}
We can take 
\[
\beta = c_1 \log(\e n/m), \quad T = \frac{c_2}{\log(\e n/m)}\delta^6.
\]
By taking $c_1$ sufficiently large and $c_2$ sufficiently small, we have, for $t \leq T$,
\[
\lambda_{n-m+1}(A_t) \geq f_{\beta, m}(A_t) \geq N_t + \frac{3}{4}.
\]
Therefore, by~\Cref{lem:maximal-inequality-for-cts-martingales}, 
\[
\Prob\{\exists t \in [0, T] : \lambda_{n-m+1}(A_t) \leq \tfrac{1}{2}\} 
\leq 
\Prob\{\exists t \in [0, T] : -N_t \geq \tfrac{1}{4}\} 
\leq \e^{-C' \delta^6/T}.
\]
If $c_2$ is sufficiently small, then 
we have 
\[
\Prob\{\exists t \in [0, T] : \lambda_{n-m+1}(A_t) \leq \tfrac{1}{2}\} 
\leq \Big(\frac{m}{\e n}\Big)^{C'/c_2} \leq \frac{m}{n},
\]
as needed. 

\subsection{Proof of~\texorpdfstring{\Cref{thm:comparison-theorem}}{the comparison theorem}}
Since $K^\circ$ is in $\Bobkov$-position, the measure $\gamma_{K^\circ}$ satisfies $\Cov(\gamma_{K^\circ}) = \alpha_{K^\circ}^2 I_n$ by~\Cref{prop:bobkov}\ref{item:scalar-covariance}. Let $X \sim \gamma_{K^\circ}$ and $Y = \alpha_{K^\circ}^{-1} X$. It holds that $\cL(Y) = \gamma_{\alpha_{K^\circ}, \alpha_{K^\circ}^{-1} K^\circ}$, which is centered, log-concave, isotropic, and supported on $\alpha_{K^\circ}^{-1} K^\circ$. Therefore, combining \Cref{prop:from-eigen-control-to-comparison-theorem,prop:eigen-stop-upper}, for any positive integer $m \leq n$ and any gauge $\gauge{\cdot}$, we obtain
\begin{align}
\E_{\gamma_n} \gauge{G} &\leq \frac{2}{\alpha_{K^\circ}^3 \sqrt{C_{\ref{prop:eigen-stop-upper}}}} \sqrt{\log(\e n/m)}\,  \E \gauge{Y} + \Lip\big(\gauge{\cdot}\big) \sqrt{2m} \nonumber  \\ 
&= \ \frac{2}{\alpha_{K^\circ}^4\sqrt{C_{\ref{prop:eigen-stop-upper}}}} \sqrt{\log(\e n/m)}\,  \E \gauge{X} + \Lip\big(\gauge{\cdot}\big) \sqrt{2m}  \nonumber \\ 
&\leq C' \Big(\sqrt{\log(\e n/m)} \E \gauge{X}  + \Lip\big(\gauge{\cdot}\big) \sqrt{m} \Big),
\label{ineq:m-final-bound-first}
\end{align}
where we used~\Cref{prop:bobkov}\ref{item:bound-on-alpha} to bound $\alpha_{K^\circ}^2$ below by $c_{\ref{prop:bobkov}}$ and we took
\[
C' = \max\Big\{\frac{2}{c_{\ref{prop:bobkov}}^2\sqrt{C_{\ref{prop:eigen-stop-upper}}}}, \sqrt{2}\Big\}. 
\]


\section{Proof of~\texorpdfstring{\Cref{thm:covering-numbers}}{the covering-number theorem}}

We will prove the following estimate in Bobkov's position. For brevity, we denote the Euclidean covering entropy of $K$ by
\[
H_K(r) = \log N(K, r B^n_2), \quad r > 0.
\]
\begin{theorem}
\label{thm:metric-entropy-bobkov}
There is a constant $C_{\ref{thm:metric-entropy-bobkov}} \in (0, \infty)$ such that for every $n \geq 1$ and every $K \in \Bodies$ with $K^\circ$ in $\Bobkov$-position and $\vol_n(K^\circ) = 1$,
\[
H_K(r) \leq C_{\ref{thm:metric-entropy-bobkov}}\begin{cases}
  n \log(1 + \tfrac{1}{r\sqrt{n}}) & r \in (0, n^{-1/2}) \\ 
 \tfrac{1}{r^2} \log(\e n r^2) & r \in [n^{-1/2}, \rad(K))  
 \end{cases}.
\]
\end{theorem}


\subsection{Proof of~\texorpdfstring{\Cref{thm:metric-entropy-bobkov}}{the metric-entropy theorem}}

The following lemma relates the covering entropy to the Gaussian width of the intersections of $K$ with Euclidean balls.

\begin{lemma}
\label{lem:one-step-bound}
There is a constant $C_{\ref{lem:one-step-bound}} \in (0, \infty)$ such that, for every $n \geq 1$, every $K \in \Bodies$, and every $r > 0$,
\[
H_K(r/2) - H_K(r) \leq C_{\ref{lem:one-step-bound}} \,
\frac{w(K \cap rB^n_2)^2}{r^2}.
\]
\end{lemma}
\begin{proof}
Let $\mathcal{P}_{r/2}$ and $\mathcal{N}_r$ denote, respectively, a maximal packing at scale $r/2$ and a minimum-cardinality covering at scale $r$. Such maximal packings are also coverings, so $|\mathcal{P}_{r/2}| \geq N(K, \tfrac{r}{2} B^n_2)$. For each $x \in \mathcal{N}_{r}$, we denote
\[
N_x = \#\{y \in \mathcal{P}_{r/2} : y \in x + rB^n_2\}.
\]
Because $\mathcal{N}_r$ is a covering at scale $r$, it clearly holds that 
\[
|\mathcal{P}_{r/2}| \leq \sum_{x \in \mathcal{N}_r} N_x.
\]
Hence, by the pigeonhole principle, there exists some $x^\star \in \mathcal{N}_r \subset K$ for which
\[
N_{x^\star} \geq \frac{|\mathcal{P}_{r/2}|}{|\mathcal{N}_r|} \geq \frac{N(K, \tfrac{r}{2} B^n_2)}{N(K, r B^n_2)}.
\]
Note that 
\[
\{y \in \mathcal{P}_{r/2} : y \in x^\star + r B^n_2\} \subset K \cap (x^\star+ rB^n_2).
\]
Moreover, the collection of points on the left-hand side above is, by definition, $r/2$-separated in the Euclidean metric. Therefore, by the Sudakov minoration,
\[
w(K \cap (x^\star + rB^n_2)) \geq c  \, r \sqrt{\log N_{x^\star}} 
\geq c\, r \, \sqrt{H_K(r/2) - H_K(r)}. 
\]
Rearranging the inequality gives
\[
H_K(r/2) - H_K(r) \leq \frac{1}{c^2} \, \Big(\frac{w(K \cap(x^\star + r B^n_2))}{r}\Big)^2.
\]
Finally, note that when $K$ is centrally symmetric, the righthand side is easily seen to be maximized over $x^\star \in K$ at the origin (\eg~\cite[Lemma 8.2]{Mou25}).
\end{proof}

The next lemma yields an estimate on these intersections when $K^\circ$ is in $\Bobkov$-position.
\begin{lemma}
\label{lem:interpolation-body-bounds}
There is a constant $C_{\ref{lem:interpolation-body-bounds}} \in (0, \infty)$ such that for every $n\geq 1$,
every $K\in \Bodies$ with $K^\circ$ in $\Bobkov$-position and $\vol_n(K^\circ) = 1$,
every $r > 0$, and every $m \in \{1, 2, \dots, n\}$,
\[
w(K \cap r B^n_2) \leq C_{\ref{lem:interpolation-body-bounds}}
\Big(\sqrt{\log(\e n/m)} + r \sqrt{m} \Big).
\]
\end{lemma}
\begin{proof}
Apply~\Cref{thm:comparison-theorem} to $\gauge{\cdot} = h_{K \cap r B^n_2}$, using $\Lip(h_{K \cap r B^n_2}) \leq r$ and the fact that \[
X \in K^\circ \subset (K \cap r B^n_2)^\circ,
\]with probability $1$.
\end{proof}


We now establish~\Cref{thm:metric-entropy-bobkov}. Let $r^\star_p = \sqrt{\tfrac{1}{p} \log(\e n/p)}$ for an integer $p \in \{1, 2, \dots, n\}$.
We will show that
\begin{equation}
\label{ineq:desired-inequality-on-integers}
H_K(r^\star_p) \leq
2 C_{\ref{lem:one-step-bound}}
C_{\ref{lem:interpolation-body-bounds}}^2
\left(
\frac{2+4c_{\ref{lem:bounded-Lipschitz-constant}}^2}{3}
+\frac{4}{9}\log 4
\right)p.
\end{equation}
If $r^\star_p \geq \rad(K)$, there is nothing to prove. Otherwise, we telescope the inequality in~\Cref{lem:one-step-bound}. Specifically, let $J \geq 1$ be an integer such that
\begin{equation}
\label{ineq:key-inequality-for-J}
2^{J-1} r^\star_p \leq \rad(K) < 2^J r^\star_p. 
\end{equation}
Then, 
\[
H_K(r^\star_p) = \sum_{j=1}^{J} \Big(H_K(2^{j-1} r^\star_p) - H_K(2^{j} r^\star_p)\Big)
\leq C_{\ref{lem:one-step-bound}} 
\sum_{j=1}^J \frac{w(K \cap 2^j r^\star_p B^n_2)^2}{4^j (r^\star_p)^2}.
\]
Applying \Cref{lem:interpolation-body-bounds}, for any $m_j \in \{1, 2 ,\dots, n\}$, we have, by our choice of $r^\star_p$,
\begin{equation}
\label{ineq:entropy-bound-almost-final}
H_K(r^\star_p) 
\leq
2 C_{\ref{lem:one-step-bound}} 
C_{\ref{lem:interpolation-body-bounds}}^2
\sum_{j=1}^J 
\Big\{m_j + \frac{p}{4^j}\frac{\log(\e n/m_j)}{\log(\e n/p)}\Big\}.
\end{equation}
We take 
\[
m_j = \ceil{\frac{p}{4^j}}, \quad \mbox{for}~1 \leq j \leq J.
\] 
Observe that $m_j \in \{1, 2, \dots, n\}$. Moreover, \eqref{ineq:key-inequality-for-J} and \Cref{lem:bounded-Lipschitz-constant} imply
\[
\frac{p}{4^J} \geq \frac{p}{4 \rad(K)^2} (r^\star_p)^2 \geq \frac{1}{4 \rad(K)^2} \geq \frac{1}{4 c_{\ref{lem:bounded-Lipschitz-constant}}^2}.
\]
Thus, for $1 \leq j \leq J$, it holds that
\[
\frac{p}{4^j} \leq m_j \leq \frac{p}{4^j} + 1 \leq \frac{p}{4^j} +
4 c_{\ref{lem:bounded-Lipschitz-constant}}^2 \frac{p}{4^J} 
\leq (1 + 4 c_{\ref{lem:bounded-Lipschitz-constant}}^2)
\frac{p}{4^j}. 
\]
Consequently, 
\begin{equation}
\label{eqn:m-sum-bound}
\sum_{j=1}^J 
m_j \leq p (1 + 4c_{\ref{lem:bounded-Lipschitz-constant}}^2)  
\sum_{j=1}^\infty 4^{-j}  =
\frac{1 + 4c_{\ref{lem:bounded-Lipschitz-constant}}^2}{3} \, p.
\end{equation}
Additionally, 
\begin{equation}
\label{eqn:log-term-bound}
\sum_{j=1}^J \frac{p}{4^j} \frac{\log(\e n/m_j)}{\log(\e n/p)} \leq 
p \sum_{j=1}^J \frac{1 + j \log 4}{4^j} = 
\Big(\frac{1}{3} + \frac{4}{9} \log 4\Big) p.  
\end{equation}
Therefore, combining the bounds~\eqref{eqn:m-sum-bound} and~\eqref{eqn:log-term-bound} with the inequality~\eqref{ineq:entropy-bound-almost-final}, 
we obtain the claimed inequality~\eqref{ineq:desired-inequality-on-integers}.

To conclude, we take the inequality~\eqref{ineq:desired-inequality-on-integers} and check that it implies the claimed entropy bound for all $r \in (0, \rad(K))$. First, fix $r > 0$; throughout we assume $r < \rad(K)$, otherwise there is nothing to prove. Then,
\[
N(K, rB^n_2) \leq N(K, r^\star_n B^n_2)N(r^\star_nB^n_2, rB^n_2) \leq
N(K, r^\star_n B^n_2) \Big(1 + 2\frac{r^\star_n}{r}\Big)^n.
\]
Taking logarithms and using $r \sqrt{n} \leq 1$ together with the estimate \eqref{ineq:desired-inequality-on-integers} at $p = n$, we obtain
\[
H_K(r) \leq H_K(r^\star_n) + n \log \Big(1 + \frac{2}{r \sqrt{n}}\Big)
\leq
\left[
\frac{
2 C_{\ref{lem:one-step-bound}}
C_{\ref{lem:interpolation-body-bounds}}^2
}{
\log 2
}
\left(
\frac{2+4c_{\ref{lem:bounded-Lipschitz-constant}}^2}{3}
+\frac{4}{9}\log 4
\right)
+2
\right]
n\log\Big(1 + \frac{1}{r \sqrt{n}}\Big).
\]
Now, for $r \geq 1/\sqrt{n}$, consider $r^\star_p$ for
\[
p = \min\{n, \ceil{q}\}, 
\quad \mbox{where} \quad 
q = \frac{\log(\e n r^2)}{r^2}.
\]
Then $r^\star_p \leq r$. Indeed, if $p = n$, then this is immediate. Otherwise, since $x \mapsto x^{-1}\log(\e n/x)$ is nonincreasing for $0 < x < \e n^2$, it holds that 
\[
(r^\star_p)^2 \leq \frac{\log(\e n/q)}{q} \leq \frac{\log(\e n r^2)}{q}  = r^2.
\]
Thus, by monotonicity of the entropy and the estimate~\eqref{ineq:desired-inequality-on-integers} applied with the above choice of $p$, we obtain
\begin{align*}
H_K(r)
&\leq
2 C_{\ref{lem:one-step-bound}}
C_{\ref{lem:interpolation-body-bounds}}^2
\left(
\frac{2+4c_{\ref{lem:bounded-Lipschitz-constant}}^2}{3}
+\frac{4}{9}\log 4
\right)
\ceil{q}
\\
&\leq
2 C_{\ref{lem:one-step-bound}}
C_{\ref{lem:interpolation-body-bounds}}^2
\left(
\frac{2+4c_{\ref{lem:bounded-Lipschitz-constant}}^2}{3}
+\frac{4}{9}\log 4
\right)
\left(1+c_{\ref{lem:bounded-Lipschitz-constant}}^2\right)
\frac{\log(\e n r^2)}{r^2},
\end{align*}
as needed. The last inequality uses the fact that, by~\Cref{lem:bounded-Lipschitz-constant},
\[
q \geq \frac{1}{r^2} \geq \frac{1}{\rad(K)^2}
\geq c_{\ref{lem:bounded-Lipschitz-constant}}^{-2}.
\]


\subsection{Proof of~\texorpdfstring{\Cref{thm:covering-numbers}}{the covering-number theorem}}

We now deduce~\Cref{thm:covering-numbers} from~\Cref{thm:metric-entropy-bobkov}.
Note the following consequence of~\Cref{thm:metric-entropy-bobkov}. 
\begin{corollary} 
\label{cor:bobkov-to-crosspolytope}
There exist universal constants $c_{\ref{cor:bobkov-to-crosspolytope}}, C_{\ref{cor:bobkov-to-crosspolytope}}>0$ such that, for every $L \in \Bodies$ with $L^\circ$ in $\Bobkov$-position and $\vol_n(L^\circ) = 1$, we have
\begin{equation}
\label{eq:bobkov-to-crosspolytope}
H_L(s)
\leq C_{\ref{cor:bobkov-to-crosspolytope}}\log N(B_1^n, c_{\ref{cor:bobkov-to-crosspolytope}}sB_2^n),
\quad \mbox{for all}~s>0.
\end{equation}
\end{corollary}
\begin{proof}
By~\Cref{lem:bounded-Lipschitz-constant},
\(\rad(L)\leq c_{\ref{lem:bounded-Lipschitz-constant}}\). Choose a sufficiently small
universal $c_{\ref{cor:bobkov-to-crosspolytope}} > 0$ so that
$c_{\ref{cor:bobkov-to-crosspolytope}}c_{\ref{lem:bounded-Lipschitz-constant}}$ lies in the range of validity of the second
estimate in~\eqref{Carsten}. If $s\geq\rad(L)$, then $H_L(s)=0$.
Otherwise, if $s < 1/\sqrt{n}$, then display \eqref{Carsten}
and~\Cref{thm:metric-entropy-bobkov} yield
\[
H_L(s)
\leq C_{\ref{thm:metric-entropy-bobkov}} \, n\log\Big(1+\frac{1}{s\sqrt{n}}\Big)
\leq C\log N(B_1^n,c_{\ref{cor:bobkov-to-crosspolytope}}sB_2^n).
\]
Suppose now that $s \geq 1/\sqrt{n}$. If $s \leq 1/(c_{\ref{cor:bobkov-to-crosspolytope}} \sqrt{n})$, then
$\log N(B_1^n,c_{\ref{cor:bobkov-to-crosspolytope}}sB_2^n)\geq cn$, and hence by 
\Cref{thm:metric-entropy-bobkov},
\[
H_L(s) \leq \frac{C_{\ref{thm:metric-entropy-bobkov}}}{s^2}\log(\e ns^2)
=C_{\ref{thm:metric-entropy-bobkov}}\,n\,\frac{\log(\e ns^2)}{ns^2}
\leq Cn \leq C''\log N(B_1^n,c_{\ref{cor:bobkov-to-crosspolytope}}sB_2^n).
\]
Finally, if $s > 1/(c_{\ref{cor:bobkov-to-crosspolytope}} \sqrt{n})$, then by the second part of display~\eqref{Carsten},
\[
\log N(B_1^n,c_{\ref{cor:bobkov-to-crosspolytope}}sB_2^n)
\geq
\frac{c}{c_{\ref{cor:bobkov-to-crosspolytope}}^2s^2}\log(1+c_{\ref{cor:bobkov-to-crosspolytope}}^2ns^2)
\geq
\frac{c'}{s^2}\log(\e ns^2) \geq c' \, H_L(s).
\]
Combining all the cases yields the claim.
\end{proof}


Now fix a convex body $K \subset \R^n$ and let
\[
K'=\frac12(K-K).
\]
Choose $S \in \GL(n)$ such that $S(K')^\circ$ is in $\Bobkov$-position and has unit volume.
By the Brunn--Minkowski and Rogers--Shephard
inequalities and the bound $\binom{2n}{n} \leq 4^n$, we have
\[
\vol_n(K)\leq \vol_n(K')
\leq 2^{-n}\binom{2n}{n}\vol_n(K)
\leq 2^n\vol_n(K).
\]
Let $T = (S^{-1})^\ast$. Then, the Blaschke--Santal\'{o} and reverse Blaschke--Santal\'{o} inequalities imply
\[
\vol_n(TK)^{1/n} \simeq 
\vol_n(TK')^{1/n} 
\simeq 
\vol_n(B^n_2)^{2/n} \simeq  \vol_n(B^n_1)^{1/n},
\]
where we used the volume normalization $\vol_n((TK')^\circ) = \vol_n(S (K')^\circ) =1$. As a result, if we set 
\[
T_\lambda = \lambda T = \Big (\frac{\vol_n(B_1^n)}{\vol_n(TK)}\Big)^{1/n} T, 
\]
then $\lambda \simeq 1$. 
Fix $x_0\in K$. Note that
\[
K-x_0\subseteq K-K=2K'.
\]
Therefore, for any $r>0$,
\begin{align*}
\log N(T_\lambda K,rB_2^n)
&=\log N\bigl(\lambda T(K-x_0),rB_2^n\bigr)\\
&\leq \log N\bigl(2\lambda TK',\tfrac r2B_2^n\bigr)\\
&=H_{TK'}\left(\frac{r}{4\lambda}\right).
\end{align*}
Applying inequality \eqref{eq:bobkov-to-crosspolytope} from~\Cref{cor:bobkov-to-crosspolytope} and using $\lambda \simeq 1$, we obtain
\[
\log N(T_\lambda K,rB_2^n)
\leq
C_{\ref{cor:bobkov-to-crosspolytope}}\log N\left(
B_1^n,\frac{c_{\ref{cor:bobkov-to-crosspolytope}}r}{4\lambda}B_2^n
\right)
\leq
C_{\ref{thm:covering-numbers}}
\log N\big(
B_1^n,c_{\ref{thm:covering-numbers}}rB_2^n
\big),
\]
as needed.

\section{Proof of~\texorpdfstring{\Cref{thm:reverse-urysohn}}{the reverse Urysohn theorem}}


For $K \in \Bodies$, we recall
\[
\rad(K) = \sup_{x \in K} \|x\|_2, \quad \text{and} \quad
\DualDvoretzky(K) = \bigg(\frac{\sqrt{n}\, M^\ast(K)}{\rad(K)}\bigg)^2.
\]
We start with the following estimate for the spherical mean width in Bobkov's position.

\begin{theorem}\label{thm:main}
There are constants $C_{\ref{thm:main}}, C_{\ref{thm:main}}' \in (0, \infty)$ such that for every $n\geq 1$
and every $K\in \Bodies$ with $K^\circ$ in $\Bobkov$-position and $\vol_n(K^\circ) = 1$,
we have
\begin{equation}\label{eq:main-bound}
\sqrt{n} \, M^\ast(K)
 \leq C_{\ref{thm:main}}\sqrt{\log\big(\e n/\DualDvoretzky(K)\big)} 
 \leq 
 C_{\ref{thm:main}}' \sqrt{\log\big(\e n\big)}.
\end{equation}
\end{theorem}

\subsection{Proof of~\texorpdfstring{\Cref{thm:main}}{the mean-width theorem}}

Let $X \sim \gamma_{K^\circ}$ and take $\gauge{\cdot} = h_K$ in~\Cref{thm:comparison-theorem}. Since $X \in K^\circ$ almost surely, we have $h_K(X)\leq 1$. Hence, integrating in polar coordinates and using $\E\|G\|_2 \geq C'' \sqrt{n}$, we find a universal constant $C > 0$ such that, for every $m \in \{1,2,\dots,n\}$,
\begin{equation}
\label{ineq:m-final-bound}
\sqrt{n} \, M^\ast(K) \leq 
C \Big(\sqrt{\log(\e n/m)}  + \rad(K) \sqrt{m} \Big).
\end{equation}
Suppose first that $\DualDvoretzky(K) \geq 4C^2$. 
Then, setting $m = \floor{\tfrac{1}{4C^2} \DualDvoretzky(K)}$ in~\eqref{ineq:m-final-bound}, we obtain
\[
\sqrt{n} \, M^\ast(K)  \leq C \sqrt{\log(8 C^2) + \log(\e n/\DualDvoretzky(K))} + 
\frac{1}{2} \sqrt{n} \, M^\ast(K). 
\]
Rearranging, we obtain 
\[
\sqrt{n} \, M^\ast(K)\leq C''' \sqrt{\log(\e n/\DualDvoretzky(K))}. 
\]
On the other hand, if $\DualDvoretzky(K) \leq 4C^2$, then $\DualDvoretzky(K) \simeq 1$, and hence we take $m = 1$ in the bound~\eqref{ineq:m-final-bound}. Using~\Cref{lem:bounded-Lipschitz-constant} to bound $\rad(K) \leq c_{\ref{lem:bounded-Lipschitz-constant}}$, we obtain 
\[
\sqrt{n} \, M^\ast(K)
\leq C \sqrt{\log (\e n)} + c_{\ref{lem:bounded-Lipschitz-constant}} 
\leq C'' \sqrt{\log(\e n/\DualDvoretzky(K))}.
\]
Combining the two cases, we obtain the first inequality. The second inequality comes from the trivial inequality $\DualDvoretzky(K) \geq c > 0$ for some universal constant $c > 0$.


\subsection{Proof of~\texorpdfstring{\Cref{thm:reverse-urysohn}}{the reverse Urysohn theorem}}

We omit the details here, as we obtain~\Cref{thm:reverse-urysohn} by passing to the difference body $K' = (K-K)/2$ and applying~\Cref{thm:main} in exactly the same way as in the proof of~\Cref{thm:quermassintegrals}; see \Cref{sec:proof-of-quermassintegrals}.



\section{Proof of~\texorpdfstring{\Cref{thm:quermassintegrals}}{the quermassintegral theorem}}

In fact, we prove the following result for the $p$th spherical mean width. It is defined as 
\[
M^\ast_p(K) = \begin{cases} 
\Big(\E_{\sigma_{n-1}} \, h_K(\theta)^p\Big)^{1/p}, & p \neq 0 \\ 
\exp \E_{\sigma_{n-1}} \log h_K(\theta), & p = 0
\end{cases}.
\]
\begin{theorem}
\label{thm:symmetric-negative-p-moments}
There exist constants $C_{\ref{thm:symmetric-negative-p-moments}}, C_{\ref{thm:symmetric-negative-p-moments}}' \in (0, \infty)$ such that,
for every $n \geq 1$ and every $K \in \Bodies$ with $K^\circ$ in $\Bobkov$-position and $\vol(K^\circ) = 1$, the following inequalities hold:
\[
M^\ast_q(K) \leq C_{\ref{thm:symmetric-negative-p-moments}} \sqrt{\frac{\log(\e n) + q}{n}} 
\quad \mbox{and} \quad 
M^\ast_{-p}(K) \leq 
C_{\ref{thm:symmetric-negative-p-moments}}' 
W_{[p]}(K) \leq 
C_{\ref{thm:symmetric-negative-p-moments}} \, 
\sqrt{\frac{\log(\e n/p)}{n}},
\]
for any $q \geq 0$ and any $p \in \{1, 2, \dots, n\}$.
\end{theorem}

\subsection{Proof of~\texorpdfstring{\Cref{thm:quermassintegrals}}{the quermassintegral theorem}}
\label{sec:proof-of-quermassintegrals} 

For a convex body $K \subset \R^n$, consider the difference body
\[
K' = \tfrac{1}{2}(K-K).
\]
Pick $S \in \GL(n)$ such that $S(K')^\circ$ has unit volume and is in $\Bobkov$-position. Set $T=(S^{-1})^\ast$ and $K''=TK'$, so that $(K'')^\circ=S(K')^\circ$.
By the second inequality in~\Cref{thm:symmetric-negative-p-moments}, for every $1 \leq k \leq n$,
\[
W_{[k]}(K'') \leq
\frac{
C_{\ref{thm:symmetric-negative-p-moments}}
}{
C_{\ref{thm:symmetric-negative-p-moments}}'
}
\sqrt{\frac{\log(\e n/k)}{n}}.
\]
For every
$E \in \Gr_{n, k}$, the Brunn--Minkowski inequality gives
\[
\vol_k(P_E TK) \leq \vol_k(P_E TK') = \vol_k(P_E K'').
\]
Therefore, integrating over $E$, we obtain
\begin{subequations}
\begin{equation}
\label{ineq:upper-on-TK}
W_{[k]}(TK) \leq W_{[k]}(K'') \leq
\frac{
C_{\ref{thm:symmetric-negative-p-moments}}
}{
C_{\ref{thm:symmetric-negative-p-moments}}'
}
\sqrt{\frac{\log(\e n/k)}{n}},
\end{equation}
for every $1 \leq k \leq n$. Finally, observe that by the Rogers--Shephard inequality,
\[
\vol_n(K'') \leq 2^{-n} \binom{2n}{n} \vol_n(TK).
\]
Therefore, the estimate $\binom{2n}{n} \leq 4^n$ and the reverse Blaschke--Santal\'o inequality imply that
\begin{equation}
\label{ineq:lower-on-volume-radius}
\vr(TK) \geq \frac{1}{2} \vr(K'') \geq \frac{c}{2} \vr((K'')^\circ)^{-1} 
= \frac{c}{2} \vol_n(B^n_2)^{1/n} \simeq \frac{1}{\sqrt{n}},
\end{equation}
\end{subequations}
as $(K'')^\circ$ was chosen to have unit volume. Above, $c > 0$ is a universal constant. 
Combining inequalities~\eqref{ineq:upper-on-TK} and~\eqref{ineq:lower-on-volume-radius} yields the claim.

\subsection{Proof of~\texorpdfstring{\Cref{thm:symmetric-negative-p-moments}}{the symmetric negative-moment theorem}}

The main estimate required to prove~\Cref{thm:symmetric-negative-p-moments} is
the following bound on the negative moments, which we prove below.
\begin{proposition}
\label{prop:main-negative-moment}
There are constants $C_{\ref{prop:main-negative-moment}}, C_{\ref{prop:main-negative-moment}}'\in (0, \infty)$ such that, for every $n\geq 1$,
every $K\in \Bodies$ with $K^\circ$ in $\Bobkov$-position and $\vol_n(K^\circ) = 1$, and
every $p \in \{1, 2, \dots, n\}$, we have
\[
\sqrt{n} \, M^\ast_{-p}(K) 
\leq 
C_{\ref{prop:main-negative-moment}}' \, 
\sqrt{n} \, W_{[p]}(K) 
\leq 
C_{\ref{prop:main-negative-moment}} \, \sqrt{\log(\e n/p)}.
\]
\end{proposition}

Assuming~\Cref{prop:main-negative-moment}, we establish~\Cref{thm:symmetric-negative-p-moments}. After possibly adjusting the universal constants, we may assume that $n$ is sufficiently large.
L\'{e}vy's concentration inequality on the sphere gives, for $q \geq 1$,
\[
\|f - \E f\|_{L^q(\sigma_{n-1})} \leq C' \Lip(f) \sqrt{\frac{q}{n}}.
\]
For $q \geq 1$, applying this to $f = h_K$ and noting that $\Lip(f) = \rad(K)$, we obtain
\begin{equation}\label{eqn:inequality-for-q-geq-1}
\sqrt{n} M^\ast_q(K) \leq \sqrt{n} \, M^\ast(K) + C' \rad(K) \sqrt{q}
\leq C'' \Big(\sqrt{\log(\e n)} + \sqrt{q}\Big), 
\end{equation}
where we used~\Cref{thm:main} and~\Cref{lem:bounded-Lipschitz-constant}.
For $q \in [0, 1]$, the monotonicity of $q \mapsto M^\ast_q(K)$ yields the result. For $p \in \{1,2, \dots, n\}$, the result follows immediately from~\Cref{prop:main-negative-moment}.


\subsection{Proof of~\texorpdfstring{\Cref{prop:main-negative-moment}}{the main negative-moment estimate}}

We recall the intrinsic volumes of a convex body $K \subset \R^n$. For $j \in \{1, \dots, n\}$, the $j$th intrinsic volume is given by
\[
V_j(K) = \binom{n}{j} \, \frac{\vol_n(B^n_2)}{\vol_j(B^j_2) \vol_{n-j}(B^{n-j}_2)} \, \E \vol_j(P_EK), 
\]
where the expectation is over a uniformly distributed $j$-dimensional subspace $E \subset \R^n$. By convention, we define $V_0(K) = 1$. 
\begin{lemma}[{\cite[Corollary 4.3 and Lemma 8.2]{Mou25}}]
\label{lem:bound-on-wills}
There is a constant ${C_{\ref{lem:bound-on-wills}} \in (0, \infty)}$ such that 
for any $n \geq 1$ and any $K \in \Bodies$,
\[
\log \Big(\sum_{j=0}^n V_j(\lambda K)\Big) 
\leq C_{\ref{lem:bound-on-wills}} \, \inf_{r > 0} \Big( \lambda \, w(K \cap r B^n_2) + \log N(K, r B^n_2) 
\Big),
\]
for any $\lambda > 0$. 
\end{lemma}

The next lemma controls the negative $p$th spherical mean width in terms of the intrinsic volumes.

\begin{lemma}
\label{lem:from-negative-moments-to-intrinsic-volumes}
There are constants $C_{\ref{lem:from-negative-moments-to-intrinsic-volumes}}, C_{\ref{lem:from-negative-moments-to-intrinsic-volumes}}'\in (0, \infty)$ such that for every $n\geq 1$
and every $K\in \Bodies$ and every integer $p \in \{1, 2, \dots, n\}$, 
\[
\sqrt{n} \, M^\ast_{-p}(K)
\leq 
C_{\ref{lem:from-negative-moments-to-intrinsic-volumes}}' 
\, 
\sqrt{n} \, W_{[p]}(K) \leq 
C_{\ref{lem:from-negative-moments-to-intrinsic-volumes}}
\, p \, V_p(K)^{1/p}.
\]
\end{lemma}
\begin{proof}
Let $L \in \Bodies$ and note that 
\[
\E_{\sigma_{n-1}} \|\theta\|_{L}^{-p} = \int_{\sphere} \rho_L(\theta)^p \, \ud \sigma_{n-1}(\theta). 
\] 
For a $p$-dimensional subspace $E$, 
\[
\frac{\vol_p(L \cap E)}{\vol_p(B^p_2)} = \int_{\mathbf{S}(E)}  \rho_L(\theta)^p \, \ud \sigma_E(\theta), 
\]
where $\mathbf{S}(E)$ denotes the sphere in $E$ and
$\sigma_E$ denotes the corresponding uniform measure.
Combining the two displays gives
\[
\E_{\sigma_{n-1}} \|\theta\|_{L}^{-p}
= \E_E  \int_{\mathbf{S}(E)} \rho_L(\theta)^p \, \ud \sigma_E(\theta) = 
\E_E \frac{\vol_p(L \cap E)}{\vol_p(B^p_2)}. 
\]
We can then apply the reverse Blaschke--Santal\'o inequality:
\[
\vol_p(L \cap E) \vol_p((L \cap E)^\circ)\geq c^p \vol_p(B^p_2)^2. 
\]
By duality of projections and sections and Jensen's inequality, this yields
\[
\E_{\sigma_{n-1}} \|\theta\|_{L}^{-p} 
\geq c^p \E \frac{\vol_p(B^p_2)}{\vol_p(P_E L^\circ)}
\geq 
c^p 
\Big(\E \frac{\vol_p(P_E L^\circ)}{\vol_p(B^p_2)}\Big)^{-1}.
\]
Set
\[
C_{n,p} = \Big(\frac{\vol_{n-p}(B^{n-p}_2)}{\vol_n(B^n_2) \binom{n}{p}}\Big)^{1/p}.
\]
Taking $L^\circ = K$, we have 
\begin{equation}
\label{ineq:upper-for-neg-p-moment}    
M^\ast_{-p}(K) \leq \frac{1}{c} W_{[p]}(K) =  \frac{1}{c} \frac{1}{\vol_p(B^p_2)^{1/p}} 
\Big(\E \vol_p(P_EK) \Big)^{1/p}
= \frac{C_{n,p}}{c} \, V_p(K)^{1/p}.
\end{equation}
Recall the formula 
\[
\vol_m(B^m_2) = \frac{\pi^{m/2}}{\Gamma(\tfrac{m}{2} + 1)}.
\]
Using $\binom{n}{p} \geq (\tfrac{n}{p})^p$ and $\Gamma(x+a) \leq (x+a)^a \Gamma(x)$ for $x, a >0$,  we have 
\[
C_{n,p} = 
\frac{1}{\sqrt{\pi}} \Big(\frac{\Gamma(\tfrac{n}{2} + 1)}{\Gamma(\tfrac{n-p}{2} + 1)}\Big)^{1/p} \binom{n}{p}^{-1/p}  \leq 
\frac{1}{\sqrt{\pi}} \frac{p}{n} \sqrt{\frac{n}{2} + 1} \leq 2\frac{p}{\sqrt{\pi n}}.
\]
Combining this with inequality~\eqref{ineq:upper-for-neg-p-moment} and then rearranging yields the claim.
\end{proof}

We are now in a position to establish~\Cref{prop:main-negative-moment}. By adjusting constants, we may assume that $n$ is large. Combining \Cref{lem:interpolation-body-bounds,lem:bound-on-wills}, we obtain
\begin{equation}
\label{ineq:general-bound}
\log \Big(\sum_{j=0}^n V_j(\lambda K)\Big) 
\leq C_{\ref{lem:bound-on-wills}}
\Big(
C_{\ref{lem:interpolation-body-bounds}}
\lambda \sqrt{\log(\e n/m)}
+C_{\ref{lem:interpolation-body-bounds}}
\lambda r \sqrt{m}
+H_K(r)
\Big),
\end{equation}
for every $\lambda, r > 0$ and $m \in \{1, \dots, n\}$. Define
\[
m^\star_p = p, \quad 
\lambda^\star_p = \frac{p}{\sqrt{\log(\e n/p)}}, 
\quad \mbox{and} \quad r^\star_p = \sqrt{\frac{\log(\e n/p)}{p}}.
\]
Assuming $r^\star_p \leq \rad(K)$ (otherwise, there is nothing to prove), \Cref{thm:metric-entropy-bobkov} gives
\begin{equation}
\label{ineq:entropy-bound}
H_K(r^\star_p) \leq C_{\ref{thm:metric-entropy-bobkov}}
\frac{p}{\log(\e n/p)} \Big\{\log\frac{\e n}{p} + \log \log \frac{\e n}{p}\Big\} 
\leq 2C_{\ref{thm:metric-entropy-bobkov}}p.
\end{equation}
Then, taking $(\lambda, r, m) = (\lambda^\star_p, r^\star_p, m^\star_p)$ in the bound~\eqref{ineq:general-bound} yields
\[
V_p(\lambda^\star_p K)
\leq \sum_{j=0}^n V_j(\lambda^\star_p K)
\leq
\exp\left(
2C_{\ref{lem:bound-on-wills}}
\big(
C_{\ref{lem:interpolation-body-bounds}}
+C_{\ref{thm:metric-entropy-bobkov}}
\big)p
\right).
\]
Since $V_p$ is positively $p$-homogeneous, taking the $p$th root and rearranging gives
\[
V_p(K)^{1/p}
\leq
\exp\left(
2C_{\ref{lem:bound-on-wills}}
\big(
C_{\ref{lem:interpolation-body-bounds}}
+C_{\ref{thm:metric-entropy-bobkov}}
\big)
\right)
\frac{\sqrt{\log(\e n/p)}}{p}.
\]
Therefore, by 
\Cref{lem:from-negative-moments-to-intrinsic-volumes}, we have 
\[
\sqrt{n} \, M^\ast_{-p}(K) \leq C_{\ref{lem:from-negative-moments-to-intrinsic-volumes}} \, p \,  V_p(K)^{1/p}  \leq
C_{\ref{lem:from-negative-moments-to-intrinsic-volumes}}
\exp\left(
2C_{\ref{lem:bound-on-wills}}
\big(
C_{\ref{lem:interpolation-body-bounds}}
+C_{\ref{thm:metric-entropy-bobkov}}
\big)
\right)
\sqrt{\log(\e n/p)},
\]
as needed. 

\bibliographystyle{abbrv}
\bibliography{references}

\end{document}
